% Источники: exam/materials/преподаватель/2020/23-03-2020-MODEL L5(2).doc; exam/materials/моделирование.pdf, разделы 12 и 32; exam/materials/прошлогодние ответы/Modelirovanie_1_-3.pdf, вопросы 12 и 32; https://github.com/HumsterProgrammer/iu9-notes/blob/main/8sem/Моделирование/Моделирование_лекция-09_17-03-26.md.
\section{Стохастическое моделирование. Проверка статистических гипотез применимости модели.}

\textbf{Статистическая гипотеза} $H$ — любое утверждение относительно функции распределения (ФР) $F(x)$ случайной величины $x$: о типе ФР, о значениях её параметров.

\textbf{Нулевая гипотеза} $H_0$ — гипотеза, справедливость которой проверяется в ходе эксперимента.
\textbf{Альтернативная гипотеза} $H_1$ — конкурирующее предположение.

Гипотеза \textbf{простая}, если её условиям удовлетворяет единственная ФР $F(x)$; иначе — \textbf{сложная}.

\subsection*{Критерий Колмогорова}

Эмпирическая функция распределения по выборке $x_1 \leqslant x_2 \leqslant \cdots \leqslant x_n$:
\[
  F_n(x) = \frac{1}{n}\sum_{i=1}^n \mathbf{1}\{x_i \leqslant x\}.
\]
Здесь $\mathbf{1}\{x_i \leqslant x\}$ равна 1, если условие выполнено, и 0 иначе; то есть числитель считает количество элементов выборки, не превосходящих $x$.

Статистика Колмогорова:
\[
  D_n = \sup_x \bigl|F_n(x) - F(x)\bigr|.
\]

Критическое значение $D_\alpha$ определяется из распределения Колмогорова при заданном уровне значимости $\alpha$. $H_0$ отвергается, если $D_n > D_\alpha$.

\subsection*{Критерий $\omega^2$ (Мизеса–Смирнова)}

Для упорядоченной выборки $x_1 \leqslant \cdots \leqslant x_n$:

\[
  \omega^2_n = \frac{1}{12n} + \sum_{i=1}^n \left[F(x_i) - \frac{2i-1}{2n}\right]^2.
\]

$H_0$ отвергается при $\omega^2_n > \omega^2_\alpha$.

\subsection*{Критерий $\chi^2$ (Пирсона)}

Применяется при большой выборке. Данные разбиваются на $k$ интервалов с наблюдаемыми частотами $\nu_i$ и теоретическими вероятностями $p_i$:
\[
  \chi^2 = \sum_{i=1}^k \frac{(\nu_i - n\,p_i)^2}{n\,p_i}.
\]
Статистика имеет распределение $\chi^2$ с $k - 1 - r$ степенями свободы ($r$ — число оценённых параметров).

\clearpage
